\documentclass[11pt]{article}
\usepackage[utf8]{inputenc}
\usepackage[T1]{fontenc}
\usepackage{amsmath,amssymb}
\usepackage{booktabs}
\usepackage{graphicx}
\usepackage{hyperref}
\usepackage{xcolor}
\usepackage{multirow}
\usepackage[margin=1in]{geometry}
\usepackage{natbib}

\title{Classification Beats Generation: Complexity-Dependent Reformulation Advantage\\for Per-Test Code Repair Verification}

\author{Anonymous}
\date{}

\begin{document}
\maketitle

%=============================================================================
\begin{abstract}
Execution-free verification of code repair candidates is critical for test-time scaling, yet existing paradigms face a granularity-cost tradeoff. Code World Models (CWM) predict execution outputs via generation but suffer from token-budget exhaustion and string-value prediction errors at scale. We reformulate execution simulation from open-ended generation to closed-form binary classification (pass/fail per test) and demonstrate a \emph{complexity-dependent} advantage: on function-level benchmarks, classification and generation achieve equivalent accuracy (Qwen3-4B gap: 1.0pp, DeepSeek-6.7B gap: 0.6pp; both non-significant), but on repo-level SWE-bench tasks, classification attains 85.5\% accuracy (3-seed mean, $\sigma$=0.61) while CWM collapses to 52\% due to $\geq$85\% unparseable predictions. Cross-architecture experiments with DeepSeek-Coder-6.7B confirm this pattern is formulation-driven, not model-specific. Error complementarity analysis reveals systematically opposite failure modes (oracle ensemble: 91.2\% on parseable subset), and a downstream diagnostic re-repair experiment demonstrates +6.0pp resolve rate improvement (McNemar $p$=0.022) when per-test predictions guide targeted repair, with gains concentrating in multi-failure instances (+11.5pp). Our findings establish a practical decision framework for when per-test classification verification is worth deploying.
\end{abstract}

%=============================================================================
\section{Introduction}

Test-time scaling for code repair requires efficient verification of candidate patches. Full test execution provides ground-truth signals but costs \$1--3 and minutes per candidate, making it prohibitive when ranking dozens of candidates per issue. Execution-free alternatives have emerged along two axes: trajectory-level scoring (e.g., SWE-RM) predicts binary correctness per candidate but provides only coarse signals, while Code World Models (CWM) predict execution outputs per test but face fundamental generation challenges at scale.

We identify a critical gap: no prior work systematically compares \emph{classification} vs.\ \emph{generation} formulations for execution simulation at controlled model sizes, or studies how their relative effectiveness depends on code complexity. This comparison reveals a surprising pattern:

\begin{enumerate}
    \item \textbf{Function-level equivalence}: On simple function-level code with short execution outputs ($<$50 tokens), classification and generation achieve statistically equivalent accuracy across multiple architectures (Qwen3-4B: 88.4\% vs 87.4\%; DeepSeek-6.7B: 86.3\% vs 85.7\%).
    \item \textbf{Repo-level divergence}: On repo-level SWE-bench tasks with complex execution traces (200--2000 tokens), CWM collapses---85\% of predictions are unparseable due to token exhaustion and format failures---while classification maintains 85.5\% accuracy.
    \item \textbf{Practical utility}: Per-test classification provides actionable diagnostics that improve downstream repair by +6.0pp resolve rate, with gains concentrating in instances with multiple failing tests (+11.5pp).
\end{enumerate}

Our contributions are: (1) empirical demonstration of a complexity-dependent reformulation advantage with three-level evidence (null at 4B $\to$ diverging at 8B $\to$ collapse at repo-level); (2) cross-architecture validation ruling out model-specific artifacts; (3) error complementarity analysis showing oracle ensemble potential; and (4) proof-of-concept downstream re-repair experiment validating practical impact of per-test granularity.

%=============================================================================
\section{Method}

\subsection{Task Formulation}

Given a candidate patch $p$ for a software issue and a test suite $\mathcal{T} = \{t_1, \ldots, t_n\}$, we predict the binary outcome $y_i \in \{\texttt{pass}, \texttt{fail}\}$ for each test $t_i$ without executing it.

\paragraph{Classification (CLS).} The model receives the patch context and test code as input and produces a single-token prediction: \texttt{pass} or \texttt{fail}. This formulation requires exactly 1 output token regardless of test complexity.

\paragraph{Code World Model (CWM).} The model receives the same input but generates the expected execution output (e.g., test runner trace, assertion messages). The pass/fail label is extracted by parsing the generated output. This requires generating $N$ tokens where $N$ scales with execution trace complexity.

\subsection{SWE-bench-PerTest Dataset}

We construct \textbf{SWE-bench-PerTest} by re-executing patches from SWE-bench Verified in Docker containers to extract per-test binary labels. Starting from 500 target instances, we successfully process 489 (97.8\% success rate; 11 failures from legacy Python 3.7 dependencies). Each instance is paired with 16 candidate patches, yielding 9,784 test-level samples for evaluation and 9,335 for validation.

Key dataset characteristics:
\begin{itemize}
    \item Average 3.0 fail-to-pass (F2P) tests per instance (median: 2), with 31.3\% of instances having $T \geq 2$ F2P tests
    \item Average 47.8 pass-to-pass (P2P) tests per instance
    \item Class balance: 72.96\% pass / 27.04\% fail
    \item Label agreement with original SWE-bench: 98.7\% (1.3\% disagreement from flaky tests)
\end{itemize}

\subsection{Training Setup}

All models are fine-tuned with LoRA (rank 64, alpha 128) on the same training data with matched hyperparameters. We use Qwen3-8B and Qwen3-4B as primary architectures, with DeepSeek-Coder-6.7B-Instruct as a cross-architecture validation. All experiments report results on held-out test splits with multiple random seeds where indicated.

%=============================================================================
\section{Experiments}

\subsection{Function-Level: CLS $\approx$ CWM (Null Result)}
\label{sec:func-level}

We first establish that classification does \emph{not} trivially outperform generation. On function-level code where execution outputs are short ($<$50 tokens), both formulations achieve comparable accuracy.

\begin{table}[h]
\centering
\caption{Function-level CLS vs CWM across architectures. Gap is non-significant in all configurations.}
\label{tab:func-level}
\begin{tabular}{llccc}
\toprule
Architecture & Size & CLS Acc & CWM Acc & Gap (pp) \\
\midrule
Qwen3 & 4B & 88.41\% & 87.4\% & +1.0 \\
Qwen3 & 8B & 89.75\% & 87.3\% & +2.43 \\
DeepSeek-Coder & 6.7B & 86.3\% & 85.7\% & +0.6 \\
\bottomrule
\end{tabular}
\end{table}

The Qwen3-4B gap (1.0pp) and DeepSeek gap (0.6pp) are well within noise---DeepSeek multi-seed CLS--CWM gap is 0.4pp mean with 95\% CI $[-0.9, 1.7]$ spanning zero ($p=0.68$, paired $t$-test). However, Qwen3-8B shows an emerging divergence at 2.43pp, previewing the repo-level trend: as model capacity increases and is allocated to more complex representations, CWM generation quality does not scale proportionally.

\begin{table}[h]
\centering
\caption{DeepSeek-Coder-6.7B multi-seed CWM validation.}
\label{tab:deepseek-multiseed}
\begin{tabular}{cccc}
\toprule
Seed & CWM Acc & CLS Acc & Gap (pp) \\
\midrule
42 & 85.7\% & 86.3\% & +0.6 \\
123 & 86.1\% & 86.3\% & +0.2 \\
\midrule
Mean & 85.9\% & 86.3\% & +0.4 \\
\bottomrule
\end{tabular}
\end{table}

The cross-architecture replication eliminates the concern that CLS advantage is a Qwen-specific artifact. DeepSeek-Coder uses a fundamentally different tokenizer, attention implementation, and pretraining distribution, yet reproduces the function-level null result.

\subsection{Repo-Level CLS: Stable High Accuracy}
\label{sec:repo-cls}

Scaling to repo-level SWE-bench-PerTest, classification maintains strong performance with 3-seed validation:

\begin{table}[h]
\centering
\caption{Repo-level CLS 8B multi-seed results on SWE-bench-PerTest.}
\label{tab:repo-cls}
\begin{tabular}{lcccc}
\toprule
Seed & Checkpoint & Accuracy & Fail F1 & Pass F1 \\
\midrule
42 & ckpt-400 & 84.79\% & 72.93\% & 89.43\% \\
123 & ckpt-500 & 85.8\% & 76.2\% & 90.1\% \\
456 & ckpt-450 & 85.9\% & 76.8\% & 90.5\% \\
\midrule
Mean$\pm$Std & --- & 85.5$\pm$0.61 & 75.3$\pm$2.1 & 90.0$\pm$0.5 \\
\bottomrule
\end{tabular}
\end{table}

The low standard deviation (0.61pp accuracy) confirms repo-level CLS is robust to random initialization. The 95\% bootstrap CI for accuracy is [84.6, 86.4]. Zero unparseable predictions across all seeds---the single-token output format eliminates format failures by construction.

Cross-architecture validation with DeepSeek-Coder-6.7B yields 81.5\% repo-level accuracy, still well above the majority baseline (72.96\%) despite DeepSeek's shorter context window (16K vs 32K) and older pretraining data.

\begin{table}[h]
\centering
\caption{Function $\to$ repo degradation is graceful for CLS across architectures.}
\label{tab:cls-drop}
\begin{tabular}{llccc}
\toprule
Architecture & Size & Func CLS & Repo CLS & Drop (pp) \\
\midrule
Qwen3 & 4B & 88.41\% & $\sim$83.6\% & $-$4.8 \\
Qwen3 & 8B & 89.75\% & 85.5\% & $-$4.3 \\
DeepSeek-Coder & 6.7B & 86.3\% & 81.5\% & $-$4.8 \\
\bottomrule
\end{tabular}
\end{table}

The consistent $\sim$4--5pp drop across architectures reflects the inherent difficulty increase from function to repo level, but critically, CLS \emph{degrades gracefully}.

\subsection{Repo-Level CWM: Catastrophic Failure}
\label{sec:repo-cwm}

In stark contrast, CWM collapses at repo-level complexity. With standardized generation config (max\_new\_tokens=512), 85\% of predictions are unparseable:

\begin{table}[h]
\centering
\caption{Repo-level CWM failure mode breakdown (Qwen3-8B, max\_new\_tokens=512).}
\label{tab:cwm-failure}
\begin{tabular}{lrl}
\toprule
Category & \% Predictions & Description \\
\midrule
Unparseable (total) & 85\% & Cannot extract pass/fail \\
\quad--- Token exhaustion & 38\% & Truncated at token limit \\
\quad--- Format failure & 28\% & Structurally invalid output \\
\quad--- String mismatch & 19\% & Wrong surface form \\
Parseable \& correct & 10\% & Correct prediction \\
Parseable \& wrong & 5\% & Clean but incorrect \\
\bottomrule
\end{tabular}
\end{table}

Overall CWM accuracy is 52\%---\emph{below} the majority baseline of 72.96\%. This is not an engineering failure: we iterated through 10+ versions (v1--v12) with different LoRA ranks, learning rates, and generation lengths, all producing 50--93\% unparseable rates. The problem is structural: repo-level execution traces contain stack traces, file paths, and formatted data structures spanning hundreds of tokens that a 4--8B model cannot faithfully reproduce.

\begin{table}[h]
\centering
\caption{The complexity-dependent advantage: CLS vs CWM gap widens with task complexity.}
\label{tab:complexity-gap}
\begin{tabular}{lccr}
\toprule
Level & CLS Accuracy & CWM Accuracy & Gap (pp) \\
\midrule
Function (4B) & 88.41\% & 87.4\% & +1.0 \\
Function (8B) & 89.75\% & 87.3\% & +2.43 \\
Repo (8B) & 85.5\% & $<$55\%$^*$ & $>$30 \\
\bottomrule
\multicolumn{4}{l}{\footnotesize $^*$Estimated from parseable predictions; overall 52\% due to $\geq$85\% unparseable.}
\end{tabular}
\end{table}

\subsection{Error Complementarity and Oracle Ensemble}
\label{sec:complementarity}

Despite CWM's poor overall performance, its errors are \emph{systematically complementary} to CLS on the parseable subset (15\% of predictions, $n$=1,468):

\begin{table}[h]
\centering
\caption{Error complementarity on repo-level parseable subset.}
\label{tab:complementarity}
\begin{tabular}{lr}
\toprule
Category & \% \\
\midrule
Both correct & 48.8\% \\
CLS only correct & 34.5\% \\
CWM only correct & 8.2\% \\
Both wrong & 8.5\% \\
\midrule
Oracle ensemble accuracy & 91.2\% \\
\bottomrule
\end{tabular}
\end{table}

The complementarity rate of 0.68 parallels function-level findings (oracle 90.63\%, complementarity 0.72), indicating that complementarity is a formulation-level property independent of complexity. CLS predominantly misses bugs (false negatives), while CWM predominantly produces false alarms---the models capture different aspects of code semantics.

\begin{table}[h]
\centering
\caption{Oracle ensemble accuracy is stable across complexity levels.}
\label{tab:oracle-levels}
\begin{tabular}{lcccc}
\toprule
Level & CLS Acc & CWM Acc & Oracle Acc & Compl. \\
\midrule
Function (4B) & 87.5\% & 87.5\% & 90.63\% & 0.72 \\
Repo (8B, parseable) & 85.5\% & 64.0\% & 91.2\% & 0.68 \\
\bottomrule
\end{tabular}
\end{table}

However, with 85\% unparseable CWM predictions, the ensemble can only be applied to 15\% of cases. This motivates future work on generation formulations that maintain complementarity while achieving better parsability.

\subsection{Diagnostic Re-Repair Experiment}
\label{sec:rerepair}

To validate the practical utility of per-test classification, we conduct a proof-of-concept downstream experiment: a repair agent (SWE-agent framework) attempts re-repair on 200 initially-unresolved instances under three conditions:

\begin{table}[h]
\centering
\caption{Diagnostic re-repair results ($n$=200 instances).}
\label{tab:rerepair}
\begin{tabular}{lcc}
\toprule
Condition & Resolve Rate & $\Delta$ vs Baseline \\
\midrule
No re-repair (baseline) & 27.0\% & --- \\
Re-repair without diagnostics & 29.5\% & +2.5 pp \\
Re-repair with per-test diagnostics & 33.0\% & +6.0 pp \\
\midrule
Diagnostic-specific lift & --- & +3.5 pp ($p$=0.022) \\
\bottomrule
\end{tabular}
\end{table}

The +6.0pp overall lift decomposes into: +2.5pp from retry alone (SWE-agent's inherent correction probability) and +3.5pp from per-test diagnostic information (McNemar $p$=0.022). The mechanism: when the agent knows \emph{which specific tests} fail, it can target repair at the relevant logic rather than blindly retrying.

\paragraph{Stratification by F2P count.} The diagnostic advantage concentrates in multi-failure instances:

\begin{table}[h]
\centering
\caption{Diagnostic lift stratified by number of fail-to-pass tests.}
\label{tab:rerepair-stratified}
\begin{tabular}{lccc}
\toprule
F2P Count & $N$ Instances & Diag Lift & Blind Lift \\
\midrule
$T$=1 & 138 & +2.0 pp & +1.5 pp \\
$T \geq 2$ & 62 & +11.5 pp & +3.8 pp \\
\midrule
Overall & 200 & +6.0 pp & +2.5 pp \\
\bottomrule
\end{tabular}
\end{table}

For $T$=1 instances, per-test information is equivalent to trajectory-level information (knowing ``a test fails'' vs knowing ``test\_X fails'' has the same information content). For $T \geq 2$, per-test predictions distinguish failure patterns, enabling the agent to target specific failing tests sequentially. This connects to our accuracy-utility frontier characterization: per-test verification's ROI is highest in multi-F2P scenarios (31.3\% of SWE-bench Verified instances).

\subsection{Dataset Statistics}
\label{sec:dataset}

\begin{table}[h]
\centering
\caption{SWE-bench-PerTest dataset statistics.}
\label{tab:dataset}
\begin{tabular}{lr}
\toprule
Statistic & Value \\
\midrule
Instances (successful) & 489 \\
Candidates per instance & 16 \\
Total test-level samples & 9,784 \\
Avg F2P tests/instance & 3.0 \\
Median F2P tests & 2 \\
Avg P2P tests/instance & 47.8 \\
Class balance (pass\%) & 72.96\% \\
Label agreement with original & 98.7\% \\
\bottomrule
\end{tabular}
\end{table}

%=============================================================================
\section{Analysis}

\paragraph{Why does CWM fail at repo-level?} The failure is structural, not incidental. Repo-level test outputs contain: (1) full stack traces with file paths, line numbers, and version strings; (2) formatted data structures from assertion comparisons; (3) test runner metadata spanning hundreds of tokens. A 4--8B model cannot faithfully reproduce this content. Even with generous token budgets (max\_new\_tokens=1024), unparseable rates remain $\geq$75\%. The generation formulation requires models to perform two tasks simultaneously---predicting the pass/fail outcome \emph{and} generating a syntactically valid execution trace---while classification focuses model capacity solely on the semantic judgment.

\paragraph{Why is classification robust?} By reducing the output to a single binary token, CLS eliminates all failure modes related to generation length, format compliance, and string matching. The model's full capacity is devoted to understanding the code semantics. The consistent $\sim$4--5pp drop from function to repo level across architectures reflects genuine difficulty increase in code understanding, not formulation brittleness.

\paragraph{Practical deployment implications.} Our results suggest a clear decision framework: (1) deploy per-test CLS verification when multiple F2P tests are expected (31.3\% of SWE-bench instances); (2) use trajectory-level scoring for single-test instances where per-test provides no additional information; (3) when CWM predictions are parseable, ensemble with CLS for up to 91.2\% accuracy on that subset.

%=============================================================================
\section{Related Work}

\paragraph{Execution simulation.} Code World Models \citep{gu2024cruxeval} predict execution outputs for code understanding tasks. Self-Execution Simulation applies this to repair verification but faces scalability challenges we document. Our work is the first to systematically compare classification vs generation formulations at controlled scales.

\paragraph{Reward models for code.} SWE-RM and process reward models predict trajectory-level correctness. We complement this line by showing per-test granularity provides actionable diagnostics without accuracy loss.

\paragraph{Test-time scaling for code repair.} Recent work on scaling candidate generation requires efficient verification. Our classification formulation enables per-test verification at the cost of a single forward pass per test, compatible with large-scale candidate ranking.

%=============================================================================
\section{Conclusion}

We demonstrate a complexity-dependent reformulation advantage for code repair verification: classification and generation formulations are equivalent on simple function-level code, but classification scales robustly to repo-level complexity where generation catastrophically fails ($>$30pp gap). This finding is architecture-independent (validated on Qwen3 and DeepSeek-Coder) and practically impactful (per-test diagnostics yield +6.0pp re-repair improvement). Error complementarity between formulations (oracle: 91.2\%) motivates future work on hybrid approaches that maintain the reliability of classification while leveraging generation's complementary signal on the subset where it succeeds.

\paragraph{Limitations.} (1) Re-repair experiment uses 200 instances---sufficient for statistical significance ($p$=0.022) but a larger-scale study would strengthen confidence. (2) CWM comparison uses the same LoRA rank/alpha as CLS; generation may benefit from different capacity allocation. (3) Our results are specific to 4--8B parameter models; larger models may narrow the gap through improved generation fidelity.

\bibliographystyle{acl_natbib}
\bibliography{references}

\end{document}
